\documentclass{article}
\usepackage{amsmath, amssymb, amsthm}
\usepackage{hyperref}
\usepackage{geometry}
\geometry{margin=1in}

\title{Erd\H{o}s Problem \#445}
\date{Last updated: 2025-08-31}
\author{Source: erdosproblems.com}

\begin{document}
\maketitle

\noindent\textbf{Status:} OPEN \quad \textbf{No prize}

\noindent\textbf{Tags:} number theory

\medskip
\noindent\textbf{Problem Statement:}

Is it true that, for any $c>1/2$, if $p$ is a sufficiently large prime then, for any $n\geq 0$, there exist $a,b\in(n,n+p^c)$ such that $ab\equiv 1\pmod{p}$?

\bigskip
\noindent\textbf{Remarks:}

Heilbronn (unpublished) proved this for $c$ sufficiently close to $1$. Heath-Brown \cite{He00} used Kloosterman sums to prove this for all $c>3/4$. 

This is discussed in this MathOverflow question.

\bigskip
\noindent\textbf{References:}

[He00] Heath-Brown, D. R., Arithmetic applications of {K}loosterman sums. Nieuw Arch. Wiskd. (5) (2000), 380--384.

\bigskip
\noindent\small{Source: \url{https://www.erdosproblems.com/445}}
\end{document}
